\section{Signs determined by the target tail}

Work in the finite-height forest space $X=J^*$. For a finite
ancestor-closed head $H$, put $P=P_H$ and $Q=I-P$; its gates $g$
are the minimal excluded nodes, and $T_g$ is the full descendant
cone rooted at $g$. We use the represented atoms, head normers,
and completed cone formula from Lemma~\ref{lem:finite-hull},
Corollary~\ref{cor:closed-atomic-ball}, and
Proposition~\ref{presentation:subtrees}.

\subsection{Signs chosen from the target tail}
\label{molb:sec:target-signs}

\begin{proposition}\label{molb:prop:target-signs}
Suppose \(Px=x\), \(Qy=y\), and \(\|x\pm y\|\le1\).  Then
\begin{equation}\label{molb:eq:target-sign-bound}
 \|y\|\le(2\sqrt2+2\sqrt6)\sqrt{1-\|x\|}
                         \le8\sqrt{1-\|x\|}.
\end{equation}
Consequently, if \(\|p\pm u_j\|\le1\) and
\(\|u_i-u_j\|\ge\varepsilon>0\) for all distinct \(i,j\) in an
infinite sequence, then
\begin{equation}\label{molb:eq:target-separated}
 \|p\|\le1-\varepsilon^2/256.
\end{equation}
\end{proposition}

\begin{proof}
The case \(x=0\) in \eqref{molb:eq:target-sign-bound} follows from
\(\|y\|\le1\).  Otherwise choose a primal unit vector \(z\)
supported on \(H\) with \(x(z)=\|x\|\).  Set
\[
 \Delta=1-\|x\|,\qquad h_g=\|y|_{T_g}\|,\qquad
 R^2=\sum_gh_g^2=\|y\|^2,\qquad\theta=1/4.
\]
If \(h_g>0\) and \(g\) has parent \(p_g\in H\), inspect the terminal
head intervals \([q,p_g]\) with absolute \(z\)-sum at least
\(\theta h_g\).  When there are any, choose the shortest such interval
and let \(d_g\in\{-1,1\}\) be the sign of its sum.  In all other cases
put \(d_g=1\).  The signs depend on \(z,y,H\), and will be the same
for both endpoint approximations.

For an atom \(A=\sum_Ia_I\mathbf1_I\), put \(b_I=z(I)\),
\(\widetilde A=\sum_Ib_I\mathbf1_I\), and
\(L_A=1-\sum_Ib_I^2\).  Let \(b_g\) be its crossing coefficient,
zero if there is no crossing.  We claim
\begin{equation}\label{molb:eq:target-charge}
 \sum_gh_g|b_g|
 \le\sum_gh_gd_gb_g+2\theta R^2+L_A/\theta.
\end{equation}
Call a gate bad if \(h_g>0\), \(d_gb_g<0\), and
\(|b_g|\ge\theta h_g\).  Its crossing head qualifies for the
selection above; it strictly contains the selected shortest suffix,
since the signs disagree.  Delete that suffix and the tail from the
crossing, retaining its nonempty initial part.  The squared-test gain
is at least \(2|b_g|\theta h_g\).  All bad gates use distinct
segments, so the deletions preserve a disjoint test family and
\[
 2\theta\sum_{g\text{ bad}}h_g|b_g|\le L_A.
\]
At remaining gates with \(d_gb_g<0\), the sum of \(h_g|b_g|\) is at
most \(\theta\sum_gh_g^2=\theta R^2\).  The identity
\(|b_g|=d_gb_g+2|b_g|\mathbf1_{\{d_gb_g<0\}}\) proves
\eqref{molb:eq:target-charge}.

Fix \(\rho>0\), and approximate each endpoint \(x+\sigma y\),
\(\sigma=\pm1\), separately within \(\rho\) by a finite atom mean.
For its expectation, the usual coefficient identity gives
\[
 \mathbb EL_A\le2(\Delta+\rho),\qquad
 \|\mathbb E\widetilde A-(x+\sigma y)\|
       \le\rho+\sqrt{2(\Delta+\rho)}=:E_\rho.
\]
Here each law has this bound because \(y(z)=0\).  If
\(r_g=\|[\mathbb E\widetilde A-(x+\sigma y)]|_{T_g}\|\), then
\(\|(r_g)\|_2\le E_\rho\).  Thus
\[
 h_g\le\mathbb E|b_g|+r_g,
 \qquad |\mathbb Eb_g-\sigma y(g)|\le r_g.
\]
Multiply by \(h_g\), sum, and apply
\eqref{molb:eq:target-charge} to obtain
\[
 R^2\le\sigma\sum_gh_gd_gy(g)+2RE_\rho
                     +2\theta R^2+2(\Delta+\rho)/\theta.
\]
All sums converge: \((h_g)\in\ell_2\), \(|y(g)|\le h_g\), and the
error-cone norms are square summable.  Both signs yield this inequality
with the same \(d_g\).  Select the sign for which the first term is
nonpositive and let \(\rho\downarrow0\).  With \(\theta=1/4\),
\[
 \tfrac12 R^2\le2R\sqrt{2\Delta}+8\Delta.
\]
Solving this quadratic proves \eqref{molb:eq:target-sign-bound}, including
\(\Delta=0\).

For the separated-sequence consequence, fix \(\rho>0\) and choose a
finite ancestor-closed \(H\) with \(\|Qp\|\le\rho\).  The head
projections of the bounded sequence have a pair with
\(\|P(u_i-u_j)\|\le2\rho\).  Let \(v=(u_i-u_j)/2\).  Convexity
gives \(\|p\pm v\|\le1\), while
\(\|Pv\|\le\rho\) and \(\|v\|\ge\varepsilon/2\).  Therefore
\[
 \|Pp\pm Qv\|\le1+2\rho,\qquad
 \|Pp\|\ge\|p\|-\rho,\qquad
 \|Qv\|\ge\varepsilon/2-\rho.
\]
Apply \eqref{molb:eq:target-sign-bound} with its constant \(8\) after
division by \(1+2\rho\), then let \(\rho\downarrow0\).  This gives
\(\varepsilon/2\le8\sqrt{1-\|p\|}\), which is
\eqref{molb:eq:target-separated}.
\end{proof}
